\newpage
\section{Podsumowanie}

W~ramach niniejszej pracy zaprojektowano i~zaimplementowano konfigurowalny framework do predykcji mutacji aminokwasowych w~białkach wirusowych z~wykorzystaniem rekurencyjnych sieci neuronowych. Opracowany system przetwarza surowe pliki FASTA poprzez pięcioetapowy potok obejmujący: czyszczenie i~filtrowanie sekwencji, podział na okresy czasowe, klasteryzację K-means z~wykorzystaniem reprezentacji ProtVec, dwukierunkowe łączenie klastrów między kolejnymi okresami oraz tworzenie zbiorów danych z~oknem przesuwnym z~ekstrakcją regionów epitopowych. Zaimplementowano trzy architektury modeli o~rosnącej złożoności --- bazowy model LSTM, model LSTM z~mechanizmem atencji Bahdanau oraz model LSTM z~podwójną atencją (DA-RNN) --- co umożliwiło systematyczne zbadanie wpływu mechanizmów atencji na jakość predykcji. Całość rozwiązania skonteneryzowano z~wykorzystaniem platformy Docker, a~konfigurację poszczególnych eksperymentów zrealizowano za pomocą plików YAML, co zapewnia reprodukowalność wyników oraz łatwość adaptacji frameworka do nowych patogenów.

% TODO: Uzupełnić po zakończeniu eksperymentów
Wyniki ewaluacji eksperymentalnej przeprowadzonej na dwóch zbiorach danych --- sekwencjach białka~S wirusa SARS-CoV-2 oraz białka HA wirusa Influenza~A podtypu H1N1 --- pozwoliły na częściową weryfikację postawionej tezy badawczej. Zastosowanie chronologicznego podziału danych, w~którym zbiór testowy obejmował wyłącznie okresy czasowe następujące po danych treningowych, zapewniło realistyczną ocenę zdolności predykcyjnych modeli, eliminując ryzyko wycieku informacji z~przyszłości do zbioru treningowego. [Placeholder: Model z~podwójną atencją osiągnął najwyższe/porównywalne wartości MCC wynoszące \ldots, co potwierdza/nie potwierdza hipotezę, że mechanizmy atencji poprawiają jakość predykcji mutacji w~regionach epitopowych.]

Przeprowadzona analiza pozwala na sformułowanie oceny zarówno zalet, jak i~ograniczeń zaproponowanego podejścia. Do istotnych zalet frameworka należy jego uniwersalność --- wykazano możliwość zastosowania tego samego potoku przetwarzania dla dwóch odrębnych wirusów przy zmianie wyłącznie pliku konfiguracyjnego YAML. Konteneryzacja Docker zapewnia pełną reprodukowalność środowiska obliczeniowego, natomiast modułowa architektura systemu ułatwia wymianę poszczególnych komponentów, na przykład zastąpienie algorytmu klasteryzacji lub modelu predykcyjnego. Wśród ograniczeń należy wskazać przede wszystkim binarny charakter predykcji --- framework przewiduje jedynie, czy w~danym regionie epitopowym nastąpi mutacja, nie precyzując jednak, jaki konkretnie aminokwas zastąpi dotychczasowy. Ponadto jakość uzyskiwanych wyników jest uzależniona od kompletności i~reprezentatywności danych dostępnych w~bazach GISAID oraz NCBI, co w~przypadku nowo pojawiających się patogenów może stanowić istotne ograniczenie. [Placeholder: Dodatkowe ograniczenia wynikające z~wyników eksperymentalnych.]

Na podstawie przeprowadzonych badań sformułowano następujące wnioski. Reprezentacje ProtVec, oparte na analogii między trójkami aminokwasów a~słowami w~przetwarzaniu języka naturalnego, stanowią skuteczny sposób kodowania sekwencji białkowych zarówno na potrzeby klasteryzacji, jak i~modelowania z~użyciem sieci neuronowych. Klasteryzacja K-means zastosowana do reprezentacji ProtVec pozwala na redukcję złożoności danych przy jednoczesnym zachowaniu informacji o~dominujących wzorcach sekwencyjnych w~poszczególnych okresach czasowych. Mechanizmy atencji umożliwiają modelowi selektywne przypisywanie różnych wag poszczególnym krokom czasowym w~oknie wejściowym, co jest szczególnie istotne w~kontekście nierównomiernej dynamiki ewolucji wirusowej. Zaproponowany framework, dzięki swojej konfigurowalności, może zostać rozszerzony na inne białka wirusowe bez konieczności modyfikacji kodu źródłowego.

Wyniki niniejszej pracy otwierają szereg kierunków dalszych badań. Po pierwsze, zasadne byłoby rozszerzenie zestawu architektur modelowych o~modele oparte na mechanizmie samoatencji (ang.~\textit{self-attention}), w~szczególności architektury Transformer, które w~wielu dziedzinach przetwarzania sekwencji wykazują przewagę nad klasycznymi modelami rekurencyjnymi. Po drugie, istotnym rozwinięciem byłoby przejście z~predykcji binarnej na predykcję konkretnych substytucji aminokwasowych, co wymagałoby przeformułowania problemu jako klasyfikacji wieloetykietowej (ang.~\textit{multi-label classification}). Po trzecie, integracja danych dotyczących trójwymiarowej struktury białek, na przykład pochodzących z~modeli AlphaFold, z~informacjami sekwencyjnymi mogłaby wzbogacić reprezentację wejściową i~poprawić jakość predykcji. Po czwarte, framework może zostać rozszerzony o~obsługę kolejnych wirusów --- na przykład HIV czy wirusa Ebola --- oraz innych białek powierzchniowych. Po piąte, systematyczna optymalizacja hiperparametrów z~wykorzystaniem metod takich jak optymalizacja bayesowska (ang.~\textit{Bayesian optimization}) mogłaby doprowadzić do poprawy wyników bez modyfikacji samej architektury modelu. Wreszcie, opracowanie modeli wielozadaniowych (ang.~\textit{multi-task learning}), łączących predykcję mutacji z~klasyfikacją wariantów wirusa, stanowi obiecujący kierunek pozwalający na jednoczesne wykorzystanie komplementarnych sygnałów zawartych w~danych.
